(1) 写$\boldsymbol v=(A,B)^T$，$\boldsymbol v'=\mathsf K\boldsymbol v$，$\mathsf K=i\begin{pmatrix}\delta&\kappa\\-\kappa&-\delta\end{pmatrix}$。矩阵满足$\mathsf K^2=(\kappa^2-\delta^2)I$，所以两个空间指数为$\pm s$，$s=\sqrt{\kappa^2-\delta^2}$，若$|\delta|>\kappa$则$s$纯虚，场为传播拍振而非指数停带。
(2) 由$\mathsf K^2=s^2I$得到传播算符$\exp(\mathsf KL)=I\cosh(sL)+\mathsf K\sinh(sL)/s$，无需逐项求本征矢。在$|\delta|<\kappa$，$s>0$，无限周期光栅存在一个随$z$增长、一个衰减的局域Bloch包络，物理有界边界条件选衰减支；停带宽度由$|\delta|=\kappa$的带缘确定。
(3) 设$A(0)=A_{\rm in}$、$B(0)=rA_{\rm in}$，在$B(L)=0$下利用传播矩阵第二行求$r=i\kappa\sinh(sL)/[s\cosh(sL)-i\delta\sinh(sL)]$，整体相位可能因后向波定义相反而改变，但反射功率唯一：$R=|r|^2=\kappa^2\sinh^2(sL)/[s^2\cosh^2(sL)+\delta^2\sinh^2(sL)]$。当$\delta=0$，$s=\kappa$，严格简化为$R=\tanh^2(\kappa L)$。
(4) 若$\kappa L\ll1$且$\delta=0$，$R\simeq(\kappa L)^2$；若$\kappa L\gg1$，$R\to1$，有限无损光栅逐渐成为高反镜。直接对方程求导有$d(|A|^2-|B|^2)/dz=0$，即前减后方向的净功率通量守恒，故$T=1-R$。出带时以$s=i\sqrt{\delta^2-\kappa^2}$解析延拓，双曲函数变三角函数，可出现透射共振。
